\section{引言：课程定位与视觉智能}

CS231N 是斯坦福大学的经典课程，全称为``Deep Learning for Computer Vision''（计算机视觉中的深度学习）。2025年春季学期是该课程创立的\textbf{十周年}，由 Fei-Fei Li、Ehsan Adeli 和 Zane Durante 共同讲授。

\begin{figure}[H]
\centering
\includegraphics[width=0.95\textwidth]{slides-images/slide-000.jpg}
\caption{CS231N 课程封面：Deep Learning for Computer Vision，Lecture 1: Introduction\protect\footnotemark}
\end{figure}
\footnotetext{Slides 第1页（Part 1）。}

\begin{figure}[H]
\centering
\includegraphics[width=0.95\textwidth]{slides-images/slide-003.jpg}
\caption{CS231N 十年历程：2015--2025 年间历届讲师阵容，包括 Andrej Karpathy、Justin Johnson、Serena Yeung 等\protect\footnotemark}
\end{figure}
\footnotetext{Slides 第4页（Part 1）。}

\subsection{AI、机器学习、深度学习与计算机视觉}

Fei-Fei Li 在开场时明确了本课程在 AI 学科图谱中的位置。AI 是一个宏大的研究领域，计算机视觉和机器学习是其中两个重要的子领域，而深度学习是机器学习的一个分支。本课程聚焦的正是\textbf{计算机视觉与深度学习的交叉地带}。

\begin{figure}[H]
\centering
\includegraphics[width=0.85\textwidth]{slides-images/slide-010.jpg}
\caption{CS231N 的学科定位：AI 大圆内部，Computer Vision 与 Deep Learning 的交集区域（红色标记），同时与 NLP、Robotics、Speech Recognition 等领域相互交叉\protect\footnotemark}
\end{figure}
\footnotetext{Slides 第11页（Part 1）。}

\begin{importantbox}{CS231N 的核心范围}
本课程不会覆盖计算机视觉或深度学习的全部内容，而是专注于二者的\textbf{核心交叉}：如何使用深度神经网络来解决视觉理解问题。同时，课程内容也与自然语言处理、机器人学等领域有密切联系。
\end{importantbox}

AI 本身是一个高度跨学科的领域，涉及数学、神经科学、计算机科学、心理学、物理学、生物学等基础学科，并在医学、法律、教育、商业等诸多应用领域产生了深远影响。Fei-Fei Li 鼓励学生将课堂所学带入各自热爱的学科方向中。

\subsection{本章小结}

CS231N 是一门聚焦深度学习与计算机视觉交叉领域的课程。AI 是一个跨学科的大领域，而计算机视觉作为智能的基石之一，与深度学习的结合催生了过去十年中最激动人心的技术革命。

%% ========== 视觉的生物学起源 ==========

